%! Author = kristofferpaulsson
%! Date = 2024-08-03

% Preamble
\documentclass[10pt,a4paper,notitlepage]{article}

\usepackage[margin=2.5cm]{geometry}
\usepackage{fontspec}
\usepackage[none]{hyphenat}  % --- NO HYPHENATION

\usepackage[hyphens]{xurl}  % --- URL HYPHENATION
\usepackage{hyperref}  % --- URL

\usepackage{draftwatermark}  % --- DONT-CITE

\urlstyle{same}  % --- URL
\pagenumbering{gobble}
\DraftwatermarkOptions{fontsize=2em,angle=0,vpos=1cm,color=black,text=Do Not Cite or Circulate}  % --- DONT-CITE

\usepackage{float}  % --- TABLE
\usepackage{makecell}  % --- TABLE
\usepackage{enumitem}  % --- LIST


\fontspec{Tinos}[
    Path = ../../font/t/,
    Extension = .ttf,
    UprightFont = *-Regular,
    Scale=1,
    Ligatures=TeX
]
\fontspec{Inter}[
    Path = ./../font/i/,
    Extension = .ttf ,
    UprightFont = *-Regular ,
    Scale=1,
    Ligatures=TeX
]

% Commands
\newcommand{\tsc}[1]{\textsc{#1}}
\newcommand*{\grc}[1]{\fontspec{Inter}#1\rmfamily}
\newcommand{\trc}[1]{\textit{\fontspec{Tinos}#1}}
\newcommand{\linkfoot}[3]{\footnote{\emph{#1}, s.v. ``{#2},'' accessed \printdate{#3}.}}

% Document
\begin{document}

    \subsection*{Ancient Greek: Reflexive Pronouns in Third Person Extra}
    \begin{itemize}[label={}, leftmargin=0]
        \item Kristoffer E.D.\ Paulsson (2026), Independent Scholar.
    \end{itemize}
    In this table, the third-person plural reflexive pronouns appear in periphrastic forms, formed by combining the
    third-person plural personal pronoun \grc{σφῶν} (sphōn) with the corresponding oblique forms of the intensive
    pronoun, yielding expressions such as \grc{σφῶν αὐτῶν} (sphōn autōn) across the oblique cases.

    It is essential to distinguish the periphrastic forms from the contracted reflexive forms, which
    are formed by crasis of \grc{ἑαυτοῦ} to \grc{αὑτοῦ}.

    \begin{table}[H]
        \begin{tabular}{c|ccc|c}
            & \multicolumn{3}{c|}{Sexus} & Genus \\
            & \tsc{Masc} & \tsc{Fem} & \tsc{Neu} & \tsc{M. F. N.} \\
            \hline
            \emph{\tsc{Plural}} & & & & \\
            \tsc{Nom} & -- & -- & -- & -- \\
            \tsc{Gen} & \makecell{\grc{σφῶν αὐτῶν} \\ \trc{sphōn autōn} \\ \small of themselves (m.)}
                    & \makecell{\grc{σφῶν αὐτῶν} \\ \trc{sphōn autōn} \\ \small of themselves (f.)}
                    & \makecell{\grc{σφῶν αὐτῶν} \\ \trc{sphōn autōn} \\ \small of themselves (n.)}
                    & \makecell{( * ) \\ \small of these items} \\
            \tsc{Dat} & \makecell{\grc{σφίσιν αὐτοῖς} \\ \trc{sphisin autois} \\ \small to/for themselves (m.)}
                    & \makecell{\grc{σφίσιν αὐταῖς} \\ \trc{sphisin autais} \\ \small to/for themselves (f.)}
                    & \makecell{--}
                    & \makecell{--} \\
            \tsc{Acc} & \makecell{\grc{σφᾶς αὐτούς} \\ \trc{sphas autous} \\ \small themselves (m.)}
                    & \makecell{\grc{σφᾶς αὐτάς} \\ \trc{sphas autas} \\ \small themselves (f.)}
                    & \makecell{--}
                    & \makecell{--} \\
        \end{tabular}
         \caption{Compiled from standard Ancient Greek grammars enhanced with transliteration and cues.
            Assume gender over sex except when the antecedent is a person.
             (Smyth, 1956:93, \S329)}
     \end{table}

    Ayer (2018) suggests that the hybrid form \grc{σφῶν ἑαυτῶν} (sphōn heautōn) is used in the plural in Ancient Greek,
    it does occur in Byzantine (Medieval) Greek.

    \subsection*{Reference List}
    \begin{itemize}[label={}, leftmargin=0]
        \item Ayer, M. ed., (2018). \textit{Goodell’s School Grammar of Attic Greek}. [online] Carlisle, Pennsylvania: Dickinson College Commentaries. Available at: \url{https://dcc.dickinson.edu/grammar/goodell/reflexive-pronouns} [Accessed 29 Mar. 2026].
        \item Smyth, H.W.\ (1956). \textit{Greek Grammar}. Cambridge, Massachusetts: Harvard University Press.
    \end{itemize}

\end{document}
